\subsection{Priprema i transformacija podataka}\label{sec:priprema}
Opisati \textbf{koji atributi su obrađeni, na koji način i zašto}. Ne morate
prikazivati Python kod, nego odluke i njihove posljedice.

\begin{table}[htbp]\centering
\caption{Primijenjeni postupci pripreme podataka}\label{tab:priprema}
\small
\begin{tabularx}{\linewidth}{@{}p{.25\linewidth}Y Y@{}}\toprule
Atribut(i) & Postupak & Razlog / posljedica \\ \midrule
\dopuni{atribut(i)} & \dopuni{npr. obrada nedostajućih vrijednosti} & \dopuni{zašto} \\
\dopuni{atribut(i)} & \dopuni{npr. kodiranje kategorija} & \dopuni{zašto} \\
\dopuni{atribut(i)} & \dopuni{npr. uklanjanje ili zadržavanje} & \dopuni{zašto} \\
\bottomrule\end{tabularx}
\end{table}

\textbf{Važna napomena.} Pripremu razdvojiti na opšte korake i transformacije
koje zavise od algoritma. Imputacija, skaliranje i kodiranje koji uče parametre
iz podataka moraju se uklopiti u odgovarajući \emph{pipeline} i učiti samo na
trening skupu. Oversampling/undersampling se ne provodi na testnom skupu.
\referencaKod{1.4 Priprema podataka}
